% Содержимое подраздела 3.2 - Фундаментальные ограничения строковых представлений

\section{Фундаментальные ограничения строковых представлений}
\label{sec:string_limitations}

Несмотря на популярность и удобство, строковые представления молекул, в первую очередь SMILES, накладывают серьезные фундаментальные ограничения на генеративные модели. Эти ограничения можно разделить на две категории: синтаксические и семантические.

Первая проблема -- \textbf{синтаксическая корректность}. Модели, работающие на символьном уровне, такие как RNN, генерируют SMILES-строку символ за символом, пытаясь выучить сложную <<грамматику>> этого языка. Однако этот процесс часто приводит к генерации синтаксически некорректных строк, которые не соответствуют ни одной реальной молекуле. Эта проблема привела к разработке подходов, которые работают не со строками, а напрямую с графовой структурой молекул. Ярким примером является \textbf{Junction Tree Variational Autoencoder (JT-VAE)} \cite{jin2018junction}. Эта модель решает проблему валидности <<по построению>>. Вместо генерации атомов, JT-VAE оперирует химически осмысленными строительными блоками (кольцами, простыми связями). Генерация происходит в два этапа: сначала создается древовидный каркас из этих блоков, а затем они собираются в итоговый молекулярный граф. Поскольку на каждом шаге используются только валидные химические компоненты, все сгенерированные молекулы гарантированно являются 100\% химически корректными.

Вторая, более глубокая проблема -- \textbf{семантическая неадекватность}. SMILES представляет молекулу на уровне <<букв>> (атомов), в то время как химики мыслят на уровне <<слов>> (функциональных групп, фрагментов). Кроме того, SMILES имеет серьезные проблемы с однозначным и инвариантным представлением хиральности -- ключевого свойства для биологической активности. Небольшие изменения в порядке обхода графа могут кардинально изменить SMILES-строку и представление стереоцентров. Эту проблему решает нотация \textbf{fragSMILES} \cite{mastrolorito2025fragsmiles}. Она представляет собой принципиальный переход от атомно-уровневого представления к фрагментному. Молекула сначала декомпозируется на химически осмысленные фрагменты, которые и становятся токенами (<<словами>>) в новой строке. Информация о хиральности кодируется явно и привязывается к точкам соединения фрагментов, что делает ее инвариантной к порядку обхода. Эксперименты показывают, что модели, обученные на fragSMILES, генерируют значительно более качественные и химически релевантные структуры, особенно в задачах, требующих точного контроля над стереохимией.
